%%%%%%%%%%%%%%%%%%%%%%%%%%%%%%%%%%%%%%%%%%%%%%%%%%%%%%%%%%%%%%%%%%
% Artigo segundo as normais mais atualizadas da ABNT
% Baseado no template de A. J. Bühler e Berg Dantas
%%%%%%%%%%%%%%%%%%%%%%%%%%%%%%%%%%%%%%%%%%%%%%%%%%%%%%%%%%%%%%%%%%
\documentclass[article,12pt,oneside,a4paper,english,brazil,sumario=tradicional]{abntex2}

% --- PACOTES ---
\usepackage[T1]{fontenc}                % Selecao de codigos de fonte.
\usepackage[utf8]{inputenc}             % Codificacao do documento
\usepackage{times}                      % Usa a fonte Times New Roman
\usepackage{indentfirst}                % Indenta o primeiro parágrafo de cada seção.
\usepackage{color}                      % Controle das cores
\usepackage{graphicx, subcaption}       % Inclusão de gráficos e subfiguras
\usepackage{microtype}                  % Para melhorias de justificação
\usepackage[abnt-emphasize=bf,abnt-and-type=e,alf,doi=true]{abntex2cite}% Citações ABNT (estilo autor-data)
\usepackage{url}
\usepackage{mathptmx}
\usepackage{amsmath}
\usepackage{hyperref}
\usepackage{multirow}				    % Permite maior formatação de tabelas
\usepackage{booktabs}                   % Para tabelas com melhor visual
\usepackage{siunitx}                    % Para formatar unidades (ex: \si{\pascal})
\usepackage{tikz}                       % Para desenhar diagramas
\usetikzlibrary{calc, positioning, arrows.meta} % Bibliotecas do TikZ
\usepackage{amsmath}                    % Para ambientes matemáticos como equation

% --- CONFIGURACOES DE PACOTES ---
\definecolor{blue}{RGB}{41,5,195}
\hypersetup{
	colorlinks=true,
	linkcolor=blue,
	citecolor=blue,
	filecolor=magenta,
	urlcolor=blue,
}

% --- CONFIGURACOES DO DOCUMENTO ---
\graphicspath{{./Figuras/}}             % Pasta para as figuras (serão usadas caixas no lugar)
\setsecheadstyle{\bfseries \normalsize \uppercase}
\setsubsecheadstyle{\normalsize \uppercase}
\setsubsubsecheadstyle{\bfseries \normalsize}
\setlrmarginsandblock{3cm}{2cm}{*}      % Margens esquerda-direita
\setulmarginsandblock{3cm}{2cm}{*}      % Margens cima-baixo
\checkandfixthelayout
\setlength{\parindent}{1.25cm}           % Tamanho do parágrafo
\OnehalfSpacing                         % Espaçamento de 1,5
\setlength{\ABNTEXcitacaorecuo}{4cm}     % Recuo para citação direta com mais de 3 linhas

% --- DADOS DO ARTIGO ---
\titulo{ESTUDO REOLÓGICO DE SUSPENSÕES DE CAULIM E VALIDAÇÃO DE UM REÔMETRO CAPILAR DE BAIXO CUSTO}
\autor{Bruno Kenji Nishitani Egami\thanks{E-mail: bruno.egami@farroupilha.ifrs.edu.br} \hspace{1em} Profª. Drª. Daniela Lupinacci Villanova\thanks{E-mail: daniela.villanova@farroupilha.ifrs.edu.br} \hspace{1em} Prof. Dr. André Zimmer\thanks{E-mail: andre.zimmer@feliz.ifrs.edu.br}}
\instituicao{
  Instituto Federal de Educação, Ciência e Tecnologia do Rio Grande do Sul - Campus Farroupilha \\
  Instituto Federal de Educação, Ciência e Tecnologia do Rio Grande do Sul - Campus Feliz \\
  Mestrado Profissional em Tecnologia e Engenharia de Materiais
}
\local{IFRS - Campus Farroupilha e Campus Feliz}
\data{\today}

\begin{document}
\selectlanguage{brazil} % Seleciona o idioma do documento
\frenchspacing          % Retira espaço extra obsoleto entre as frases.

% --- CABEÇALHO ---
\maketitle % Gera o cabeçalho automaticamente com os dados definidos acima

\begin{resumo}
\noindent
Este estudo avalia as características de fluxo de suspensões aquosas de caulim em diversas concentrações e apresenta a validação de um reômetro capilar acessível para tal caracterização. O controle reológico é fundamental no processamento cerâmico, impactando a aplicabilidade de técnicas avançadas como a impressão 3D. A formulação de pastas cerâmicas com elevada fração de sólidos é um objetivo tecnológico, frequentemente dificultado pela plasticidade inerentemente baixa da matéria-prima, necessitando a adição de modificadores reológicos. Suspensões foram formuladas com concentrações de caulim variando entre 60\% e 75\% em massa. As propriedades reológicas foram mensuradas utilizando um reômetro comercial (Anton Paar MCR 102) como padrão de referência e comparadas com as medições obtidas no reômetro capilar proposto. Os dados de ambos os dispositivos mostraram boa correlação, especialmente em taxas de deformação elevadas, confirmando a adequação do equipamento de baixo custo. As suspensões com 60\% a 65\% de sólidos, que se mostraram passíveis de teste, exibiram comportamento de afinamento por cisalhamento (\textit{shear-thinning}). O trabalho demonstra o aumento significativo da tensão de escoamento conforme a concentração de sólidos se eleva, fornecendo uma base experimental para o desenvolvimento futuro de pastas cerâmicas aditivadas destinadas à manufatura aditiva.
\vspace{\onelineskip}
\noindent
\textbf{Palavras-chave}: Caulim. Reologia. Suspensões Cerâmicas. Reômetro Capilar. Baixo Custo. Impressão 3D.
\end{resumo}

\textual
\pagestyle{simple}
\aliaspagestyle{chapter}{simple}
\thispagestyle{empty}

% --- SEÇÕES DO ARTIGO ---

\section{Introdução}
\label{sec:introducao}
\normalsize
% Reescrevendo a Introdução para evitar autoplágio
A relação entre estrutura, processamento e propriedades é um pilar da ciência e engenharia de materiais, sendo particularly crítica para materiais cerâmicos. Processos de fabricação cerâmica que utilizam via úmida, tais como colagem de barbotina, extrusão e atomização, são fortemente dependentes do controle das características de fluxo das suspensões empregadas. Este controle é o objeto de estudo da reologia, que se torna uma ferramenta indispensável no desenvolvimento de suspensões que conciliem alta concentração de sólidos com fluidez adequada. Tal combinação é necessária para a obtenção de produtos finais com alta densidade e, consequentemente, propriedades mecânicas otimizadas, visto que a redução da porosidade (baixa absorção de água) está correlacionada com melhor desempenho mecânico \citeonline{rocha2007}.

Nos últimos anos, a manufatura aditiva (AM), comumente chamada de impressão 3D, consolidou-se como uma rota tecnológica promissora para a fabricação de componentes cerâmicos. Entre as diversas técnicas de AM, a \textit{Direct Ink Writing} (DIW), ou \textit{Robocasting}, ganha destaque. Este método baseia-se na extrusão controlada de pastas altamente viscosas, camada por camada, para construir objetos tridimensionais complexos \citeonline{shahzad2021, dewitte2022}. O processo exige pastas com elevada fração volumétrica de sólidos e um comportamento reológico precisamente ajustado para garantir a integridade estrutural das camadas depositadas.

A impressão 3D de cerâmicas, exemplificada na Figura \ref{fig:exemplos_impressao}, proporciona uma liberdade geométrica sem precedentes, permitindo a fabricação de peças com geometrias intrincadas, muitas vezes irrealizáveis por métodos convencionais \citeonline{oliveira2024, zocca2015}. Contudo, o sucesso da técnica DIW está intrinsecamente ligado a um controle reológico preciso. \citeonline{ordonez2019} e \citeonline{pires2023} enfatizam que as pastas para DIW devem exibir propriedades específicas, preferencialmente um comportamento pseudoplástico (afinamento por cisalhamento) ou de fluido de Bingham. É essencial que a pasta flua sob a tensão aplicada durante a extrusão, mas que também recupere rapidamente sua viscosidade após a deposição (tixotropia), garantindo a capacidade de suportar o peso das camadas subsequentes sem colapsar (propriedade conhecida como \textit{buildability}).

\begin{figure}[htbp]
  \centering
  % Placeholder mantido como solicitado
  \framebox{\parbox{0.8\textwidth}{\centering
    \vspace{3cm}
    \textbf{Figura 1 Placeholder} \\
    \small\textit{Exemplos de peças em argila confeccionadas com impressão 3D.}
    \vspace{3cm}
  }}
  \caption{Exemplos de peças em argila confeccionadas com impressão 3D.}
  \fonte{(a) e (b) \citeonline{3dpotter2025}; (c) \citeonline{knapp2025}.}
  \label{fig:exemplos_impressao}
\end{figure}

Neste cenário, a caracterização detalhada das propriedades reológicas da matéria-prima base torna-se um pré-requisito crucial. Entretanto, muitas argilas cauliníticas, como o metacaulim, demonstram baixa plasticidade intrínseca, o que dificulta ou inviabiliza a impressão direta sem o uso de aditivos modificadores. Pesquisas como as de \citeonline{campos2023} e \citeonline{revelo2019} exploram o uso de aditivos, como bentonita ou cinza volante, para aprimorar a plasticidade, a trabalhabilidade e a qualidade superficial de pastas de caulim destinadas à DIW. Uma estratégia alternativa, investigada por \citeonline{ordonez2019}, concentra-se no emprego de defloculantes (como poliacrilato de sódio) para formular pastas fluidas com teores de sólidos excepcionalmente altos.

Este trabalho se insere como uma investigação preliminar, visando futuras aplicações em impressão 3D. Seus objetivos são duplos: (1) investigar o comportamento reológico fundamental de suspensões aquosas de caulim em várias concentrações e (2) construir e validar um dispositivo de reometria capilar de baixo custo como alternativa acessível aos equipamentos comerciais. A validação de instrumentação de baixo custo representa, por si só, um avanço importante para democratizar a pesquisa e o desenvolvimento tecnológico, como já demonstrado por \citeonline{morita2005} em um trabalho similar de validação de um módulo de reometria capilar de baixo custo integrado a uma máquina de ensaios universal.

\section{Fundamentação Teórica}
\label{sec:teoria}

\subsection{Princípios Fundamentais e Modelagem}
% Texto mantido, pois a reescrita poderia comprometer a precisão técnica das definições.
O comportamento de um fluido sob cisalhamento é descrito pela relação entre a \textbf{tensão de cisalhamento} ($\tau$) e a \textbf{taxa de cisalhamento} ($\dot{\gamma}$). A \textbf{viscosidade} ($\eta$) é a propriedade que relaciona esses dois parâmetros, definida pela Equação \ref{eq:viscosidade}:
\begin{equation} \label{eq:viscosidade}
    \eta = \frac{\tau}{\dot{\gamma}}
\end{equation}
Suspensões cerâmicas, por serem sistemas multifásicos, raramente se comportam como fluidos Newtonianos (onde a viscosidade é constante), exibindo um comportamento \textbf{não-Newtoniano} complexo \citeonline{rocha2007}.

Para descrever matematicamente esse comportamento, diversos modelos reológicos podem ser empregados. O modelo de \textbf{Herschel-Bulkley} é comumente utilizado para suspensões concentradas, pois engloba três parâmetros físicos importantes (Equação \ref{eq:herschel-bulkley}):
\begin{equation} \label{eq:herschel-bulkley}
    \tau = \tau_y + K\dot{\gamma}^n
\end{equation}
Nesta equação, $\tau_y$ é a \textbf{tensão de escoamento} (ou \textit{yield stress}), que representa a tensão mínima necessária para iniciar o fluxo, vencendo a rede de interações interpartículas em repouso. O parâmetro $K$ é o \textbf{índice de consistência}, relacionado à magnitude da viscosidade do material. Por fim, $n$ é o \textbf{índice de comportamento de fluxo}, que classifica o fluido.

Para a maioria das suspensões cerâmicas, observa-se o comportamento \textbf{pseudoplástico} ou \textit{shear-thinning}, onde $n < 1$. Este comportamento é caracterizado pela diminuição da viscosidade com o aumento da taxa de cisalhamento. O fenômeno é atribuído à quebra de aglomerados e, no caso do caulim, à orientação das partículas anisometricas (lamelares) na direção do fluxo. Este é o comportamento desejado para DIW, pois permite que o material flua facilmente pelo bico (alta taxa de cisalhamento), mas mantenha sua forma após a deposição (baixa taxa de cisalhamento).

Adicionalmente, muitas suspensões exibem \textbf{tixotropia}, um fenômeno dependente do tempo no qual a viscosidade diminui sob cisalhamento constante e se recupera gradualmente quando o cisalhamento cessa. A tixotropia está associada à quebra e reconstrução da estrutura interna da suspensão e é uma característica crucial para a impressão 3D.

\subsection{O Papel dos Aditivos na Reologia de Pastas Cerâmicas}
% Texto reescrito para diversificar a linguagem
O processamento de suspensões puras de caulim e água, foco deste estudo, enfrenta uma limitação intrínseca: o aumento da fração de sólidos resulta em um crescimento exponencial da tensão de escoamento e da viscosidade. Isso pode tornar o material excessivamente rígido, não plástico e inadequado para extrusão. Para contornar essa dificuldade e possibilitar a formulação de pastas com alta concentração de sólidos (>65-70\%), adequadas para DIW \citeonline{revelo2018}, a utilização de aditivos torna-se essencial.
\begin{itemize}
    \item \textbf{Plastificantes:} Algumas argilas, como o metacaulim, podem exibir baixa plasticidade. Aditivos que conferem plasticidade, como a \textbf{bentonita} (uma argila do grupo das esmectitas com alta capacidade de absorção de água), são frequentemente incorporados em pequenas proporções (tipicamente 4-5\%). A bentonita atua como agente ligante e espessante, melhorando a coesão, a trabalhabilidade e a resistência mecânica da pasta no estado "verde" (antes da queima), viabilizando sua impressão.
    \item \textbf{Defloculantes:} Para alcançar teores de sólidos ainda mais elevados, empregam-se defloculantes (ou dispersantes). Compostos como \textbf{silicato de sódio} ou \textbf{poliacrilato de sódio} modificam as cargas elétricas na superfície das partículas de argila, induzindo repulsão eletrostática entre elas. Esse efeito impede a formação de aglomerados (flocos), libera a água aprisionada nessas estruturas e causa uma redução acentuada na viscosidade da suspensão. Isso permite a criação de pastas fluidas mesmo com um teor reduzido de água.
    \item \textbf{Aditivos Minerais (Cargas):** Outros pós minerais podem ser adicionados para otimizar as propriedades da pasta. \citeonline{revelo2019} observaram que a incorporação de \textbf{cinza volante} (\textit{fly ash}) resultou em melhorias notáveis na trabalhabilidade e no acabamento superficial de pastas de caulim. Atribui-se esse efeito à morfologia esférica das partículas de cinza volante, que agem como micro-rolamentos, diminuindo o atrito interno e facilitando o fluxo da mistura.
\end{itemize}

\section{Metodologia Experimental}
\label{sec:metodologia}

\subsection{Materiais e Preparação das Amostras}
% Texto reescrito
Utilizou-se caulim de grau industrial, fornecido pela empresa BRX Minérios (São Caetano do Sul/SP), para os experimentos. A preparação das suspensões consistiu na pesagem das quantidades predefinidas de caulim e água destilada, que foram então transferidas para um saco plástico resistente. A homogeneização foi realizada por meio de agitação manual vigorosa até a obtenção de uma pasta visualmente uniforme e sem grumos. A utilização de sacos plásticos também simplificou o processo de alimentação do reômetro: uma das extremidades do saco era cortada, permitindo a transferência controlada da pasta para o reservatório do equipamento. As composições exatas das suspensões, expressas em porcentagem mássica de caulim e água, estão listadas na Tabela \ref{tab:concentracoes}.

\begin{table}[htbp]
\centering
\caption{Formulações das suspensões de caulim analisadas.}
\label{tab:concentracoes}
\begin{tabular}{@{}ccc@{}}
\toprule
\textbf{Amostra} & \textbf{\% de Caulim (Sólidos)} & \textbf{\% de Água} \\ \midrule
A & 75,0\% & 25,0\% \\
B & 70,0\% & 30,0\% \\
C & 67,5\% & 32,5\% \\
D & 65,0\% & 35,0\% \\
E & 62,5\% & 37,5\% \\
F & 60,0\% & 40,0\% \\ \bottomrule
\end{tabular}
\end{table}

Durante a etapa de preparação, foi constatado que as formulações contendo 32,5\% ou mais de água (Amostras C a F) resultaram em pastas com consistência adequada e boa trabalhabilidade manual. Em contrapartida, as formulações com menor teor de água (Amostras A e B, correspondendo a 70\% e 75\% de sólidos, respectivamente) não foram incluídas nos ensaios reológicos subsequentes. Estas amostras exibiram plasticidade insuficiente, resultando em uma massa quebradiça e sem coesão ("aspecto de farofa"), que se desintegrava facilmente sob pressão manual, inviabilizando sua conformação e teste. Essa observação prática ressalta a importância de um teor mínimo de água para conferir plasticidade ao caulim puro, uma dificuldade também encontrada por \citeonline{campos2023} ao trabalhar com metacaulim.

\subsection{Equipamentos de Análise Reológica}
% Texto reescrito
A caracterização reológica das suspensões foi realizada utilizando dois equipamentos distintos, que são detalhados a seguir e visualizados na Figura \ref{fig:reometros}.
\begin{itemize}
    \item \textbf{Reômetro Rotacional Comercial:} Como equipamento de referência, utilizou-se um reômetro modelo MCR 102 da marca Anton Paar. Este equipamento, dotado de geometria de placas paralelas com 25 mm de diâmetro, é especialmente adequado para medições precisas em baixas taxas de cisalhamento. A amostra F (60\% de sólidos) foi caracterizada neste dispositivo para fins comparativos.
    \item \textbf{Reômetro Capilar de Baixo Custo:} Foi projetado e construído um reômetro do tipo capilar, empregando componentes de fácil aquisição e peças fabricadas por impressão 3D. O objetivo principal foi desenvolver uma alternativa de custo reduzido para análises reológicas, diminuindo a dependência de equipamentos comerciais de alto valor.
\end{itemize}

\begin{figure}[htbp]
  \centering
  \begin{subfigure}[b]{0.45\textwidth}
      \centering
      % Placeholder mantido
      \framebox{\parbox{0.8\linewidth}{\centering
        \vspace{3cm}
        \textbf{Figura 2a Placeholder} \\
        \small\textit{Reômetro rotacional comercial.}
        \vspace{3cm}
      }}
      \caption{Reômetro rotacional comercial.}
      \label{fig:reometro_comercial}
  \end{subfigure}
  \hfill
  \begin{subfigure}[b]{0.45\textwidth}
      \centering
      % Placeholder mantido
      \framebox{\parbox{0.8\linewidth}{\centering
        \vspace{3cm}
        \textbf{Figura 2b Placeholder} \\
        \small\textit{Reômetro capilar de baixo custo.}
        \vspace{3cm}
      }}
      \caption{Reômetro capilar de baixo custo.}
      \label{fig:reometro_diy}
  \end{subfigure}
  \caption{Equipamentos utilizados para a análise reológica.}
  \label{fig:reometros}
\end{figure}

\subsection{Desenvolvimento do Reômetro Capilar de Baixo Custo}
\label{subsubsec:reometro-capilar}
% Texto significativamente reescrito para evitar autoplágio
O sistema de reometria capilar desenvolvido integra componentes de hardware e software. O hardware é constituído por um corpo principal, um pistão e um conjunto de capilares intercambiáveis, todos fabricados com tecnologia de impressão 3D (polímero não especificado, mas provavelmente ABS ou PLA). A aquisição de dados é realizada por um sistema eletrônico baseado em um microcontrolador Arduino Pro Micro, que interpreta o sinal analógico proveniente de um transdutor de pressão com faixa de medição de até \SI{150}{psi}. O firmware embarcado no Arduino executa a conversão do sinal do sensor para valores de tensão e os transmite para um computador via interface serial. A Figura \ref{fig:esquema_reometro} apresenta um diagrama esquemático do arranjo experimental.

Uma distinção importante em relação ao dispositivo de baixo custo descrito por \citeonline{morita2005} é o método de aplicação da força de extrusão. Enquanto o sistema de Morita et al. foi concebido para ser acoplado a uma Máquina Universal de Ensaios, que impõe o deslocamento do pistão a uma velocidade controlada, o reômetro deste trabalho emprega um sistema pneumático (compressor de ar) para gerar a pressão de extrusão, a qual é medida diretamente pelo transdutor de pressão.

A interface com o usuário, o controle da sequência experimental e o tratamento dos dados são gerenciados por um conjunto de scripts desenvolvidos na linguagem Python. O script principal, denominado \texttt{1.Controle\_Reometro.py}, oferece uma interface para o operador, executa a calibração do sensor de pressão e controla o processo de aquisição de dados. Uma funcionalidade relevante implementada neste script é um sistema de gatilho (trigger) baseado em níveis de pressão pré-definidos. Esse sistema aciona e interrompe automaticamente um cronômetro digital, permitindo uma medição mais precisa do tempo necessário para a extrusão de um volume conhecido de material. Após a extrusão em cada nível de pressão, o operador é solicitado a inserir a massa do material extrudado, medida externamente em uma balança de precisão. Todos os dados coletados (pressão, massa, tempo, dimensões do capilar) e metadados relevantes são armazenados em arquivos no formato JSON.

Um segundo script, \texttt{2.Analise\_reologica.py}, é responsável pelo pós-processamento completo dos dados brutos. Ele realiza a leitura dos arquivos JSON, calcula os parâmetros reológicos aparentes (tensão e taxa de cisalhamento na parede), aplica as correções reológicas necessárias (discutidas na Seção \ref{subsec:correcoes}) e realiza o ajuste dos dados corrigidos a cinco modelos reológicos distintos (Newtoniano, Lei da Potência, Bingham, Herschel-Bulkley e Casson). O script seleciona automaticamente o modelo que melhor descreve o comportamento do fluido, com base no maior coeficiente de determinação ($R^2$). Scripts adicionais foram desenvolvidos para facilitar a visualização gráfica dos resultados e a comparação entre diferentes amostras ou condições experimentais. A documentação técnica detalhada do projeto, incluindo esquemas, códigos-fonte e arquivos de modelos 3D, está publicamente disponível no repositório online referenciado em \citeonline{egami2025}.

\begin{figure}[htbp]
\centering
\resizebox{0.95\textwidth}{!}{%
% Código TikZ mantido
\begin{tikzpicture}[node distance=2cm, auto]
    % Carregar bibliotecas
    \usetikzlibrary{calc, positioning}

    % Definindo os estilos
    \tikzstyle{block} = [rectangle, draw, fill=blue!20, text width=8em, text centered, rounded corners, minimum height=3em]
    \tikzstyle{instrument} = [circle, draw, fill=orange!30, text centered, minimum size=5em, align=center]
    \tikzstyle{data_block} = [rectangle, draw, fill=purple!20, text width=8em, text centered, rounded corners, minimum height=3em]
    \tikzstyle{line} = [draw, thick, -{Stealth[length=3mm]}]
    \tikzstyle{data_line} = [draw, thick, dashed, color=green!60!black, -{Stealth[length=3mm]}]
    \tikzstyle{manual_line} = [draw, thick, dotted, color=red!80!black, -{Stealth[length=3mm]}]

    % Posicionando os nós (componentes) na horizontal
    \node [block] (compressor) {Compressor};
    \node [block, right=1.5cm of compressor] (regulator) {Registro Regulador de Pressão};
    \node [instrument, right=1.5cm of regulator] (manometer) {Manômetro};
    \node [block, right=1.5cm of manometer] (valve) {Válvula de Bloqueio};
    \node [block, right=1.5cm of valve] (rheometer) {Corpo do Reômetro};

    % Ponto de tomada de pressão (entre a válvula e o reômetro)
    \coordinate (tap_point) at ($(valve.east)!0.5!(rheometer.west)$);

    % Linha de Dados (abaixo da linha principal)
    \node [instrument, below=2cm of tap_point] (transducer) {Transdutor \\ de Pressão};
    \node [data_block, left=2cm of transducer] (arduino) {Arduino};
    \node [data_block, left=2cm of arduino] (computer) {Computador (Python)};

    % Balança para medição de massa, posicionada mais abaixo do reômetro para evitar sobreposição
    \node [data_block, below=3.5cm of rheometer] (balance) {Balança};

    % Desenhando as linhas de conexão
    % Linha Pneumática Principal
    \path [line] (compressor) -- (regulator);
    \path [line] (regulator) -- (manometer);
    \path [line] (manometer) -- (valve);
    \path [line] (valve) -- (rheometer);

    % Linha de Pressão para o Transdutor
    \path [line] (tap_point) -- (transducer);

    % Linha de Dados Automática
    \path [data_line] (transducer) -- (arduino);
    \path [data_line] (arduino) -- (computer);

    % Fluxo de material para a balança (agora como coleta manual)
    \path [manual_line] (rheometer) -- (balance);

    % Linha de Dados Manual em "L"
    \path [manual_line] (balance) -| (computer.south);
\end{tikzpicture}
}
\caption{Esquema de funcionamento do reômetro capilar, mostrando a linha pneumática (sólida), a linha de aquisição de dados automática (tracejada) e as linhas de coleta e entrada de dados manual (pontilhada).}
\label{fig:esquema_reometro}
\end{figure}


\subsection{Cálculos Reológicos do Reômetro Capilar}
% Texto mantido, pois a reescrita poderia comprometer a precisão das equações e definições.
Os parâmetros fundamentais (tensão de cisalhamento, $\tau$; taxa de cisalhamento, $\dot{\gamma}$; e viscosidade, $\eta$) são obtidos indiretamente. A \textbf{tensão de cisalhamento na parede} ($\tau_w$) é calculada a partir da pressão de extrusão ($P$) e da geometria do capilar (Raio $R$, Comprimento $L$), conforme \citeonline{barnes1989} (Equação \ref{eq:tau_w}):
\begin{equation} \label{eq:tau_w}
    \tau_w = \frac{P \cdot R}{2 \cdot L}
\end{equation}

A \textbf{taxa de cisalhamento aparente na parede} ($\dot{\gamma}_{aw}$) é calculada a partir da vazão volumétrica ($Q$), obtida pela massa extrudada ($m$) no tempo ($t$) e a densidade da pasta ($\rho$). A equação (\citeonline{reed1995}) é (Equação \ref{eq:gamma_aw}):
\begin{equation} \label{eq:gamma_aw}
    \dot{\gamma}_{aw} = \frac{4 \cdot Q}{\pi \cdot R^3}
\end{equation}
Este valor é "aparente" pois assume um perfil de velocidade parabólico, válido apenas para fluidos Newtonianos. A \textbf{viscosidade aparente} ($\eta_a$) é a razão entre esses dois valores (Equação \ref{eq:eta_a}):
\begin{equation} \label{eq:eta_a}
    \eta_a = \frac{\tau_w}{\dot{\gamma}_{aw}}
\end{equation}

\subsection{Correções para Fluidos Não-Newtonianos}
\label{subsec:correcoes}
% Texto mantido para precisão das definições.
A reometria capilar para fluidos não-Newtonianos requer a aplicação de correções matemáticas para compensar perfis de velocidade não parabólicos e outros efeitos, conforme detalhado por \citeonline{rocha2007} e \citeonline{barnes1989}.

\subsubsection{Weissenberg-Rabinowitsch}
Esta é a correção mais importante para fluidos pseudoplásticos. Ela converte a taxa de cisalhamento aparente ($\dot{\gamma}_{aw}$) na taxa de cisalhamento verdadeira na parede ($\dot{\gamma}_w$), ajustando para o perfil de velocidade "achatado" desses fluidos. A equação (\citeonline{barnes1989}) é (Equação \ref{eq:wr}):
\begin{equation} \label{eq:wr}
    \dot{\gamma}_w = \frac{\dot{\gamma}_{aw}}{4} \left( 3 + \frac{d(\ln \dot{\gamma}_{aw})}{d(\ln \tau_w)} \right)
\end{equation}
O termo diferencial na equação, $d(\ln \dot{\gamma}_{aw})/d(\ln \tau_w)$, representa o inverso da inclinação da curva de fluxo em um gráfico log-log de $\tau_w$ versus $\dot{\gamma}_{aw}$. Este termo, frequentemente denotado como $1/n'$, onde $n'$ é o índice de comportamento de fluxo local, é calculado numericamente a partir dos dados experimentais.

\subsubsection{Bagley}
Esta correção quantifica as perdas de pressão $(P_e)$ que ocorrem na entrada e saída do capilar, que não são devidas ao cisalhamento viscoso. A sua determinação exige ensaios com capilares de diferentes comprimentos ($L$), conforme a Equação \ref{eq:bagley}:
\begin{equation} \label{eq:bagley}
    P = 2 \tau_w \left( \frac{L}{R} \right) + P_e
\end{equation}

\subsubsection{Mooney}
Esta correção avalia o fenômeno de \textbf{deslizamento na parede} (\textit{wall slip}), comum em suspensões concentradas, onde o fluido escorrega na parede do capilar, resultando em uma viscosidade artificialmente baixa. Requer ensaios com capilares de raios ($R$) distintos (Equação \ref{eq:mooney_model}):
\begin{equation} \label{eq:mooney_model}
    \dot{\gamma}_{aw} = \dot{\gamma}_{s,f} + \frac{C}{R}
\end{equation}

\subsection{Procedimento Experimental}
% Texto reescrito
Os experimentos com o reômetro capilar foram conduzidos utilizando um único capilar com dimensões fixas: diâmetro interno de \SI{2}{\milli\meter} e comprimento de \SI{30}{\milli\meter}. As suspensões analisadas foram as correspondentes às amostras D, E e F, que possuem teores de sólidos de 65\%, 62,5\% e 60\% em massa, respectivamente.

Nesta investigação preliminar, optou-se por aplicar apenas a correção de Weissenberg-Rabinowitsch aos dados de taxa de cisalhamento aparente. É importante reconhecer que a não aplicação das correções de Bagley (efeitos de entrada/saída) e Mooney (deslizamento na parede), diferentemente da abordagem adotada no estudo de validação de \citeonline{morita2005}, pode levar a um certo desvio nos valores absolutos calculados para a tensão de cisalhamento e a viscosidade. No entanto, considerando que o objetivo primário deste trabalho era a comparação das tendências gerais de comportamento reológico entre as amostras e a validação da concordância entre o reômetro de baixo custo e o equipamento comercial, esta simplificação foi considerada metodologicamente aceitável para esta fase do estudo.

\section{Resultados e Discussão}
\label{sec:resultados}
% Texto reescrito
As Figuras \ref{fig:fluxo_comparativo} e \ref{fig:viscosidade_comparativa} apresentam os resultados comparativos obtidos nos ensaios reológicos.

\begin{figure}[htbp]
    \centering
    % Placeholder mantido
    \framebox{\parbox{0.8\textwidth}{\centering
        \vspace{3cm}
        \textbf{Figura 4 Placeholder} \\
        \small\textit{Comparativo das curvas de fluxo (Tensão de Cisalhamento vs. Taxa de Cisalhamento) em escala logarítmica.}
        \vspace{3cm}
      }}
    \caption{Comparativo das curvas de fluxo (Tensão de Cisalhamento vs. Taxa de Cisalhamento) em escala logarítmica, obtidas no reômetro comercial (MCR102) e no reômetro capilar de baixo custo para diferentes concentrações.}
    \label{fig:fluxo_comparativo}
\end{figure}

A inspeção dos gráficos indica que todas as amostras testadas exibem um comportamento reológico não-newtoniano do tipo pseudoplástico, também conhecido como \textit{shear-thinning}. Este comportamento é evidenciado pela redução da viscosidade aparente conforme a taxa de cisalhamento aumenta, sendo uma característica bem documentada para suspensões aquosas de caulim na literatura científica \citeonline[e.g.,][]{rocha2007, revelo2018}.

O achado central desta seção reside na validação do reômetro capilar de baixo custo. A comparação dos dados obtidos para a Amostra F (60\% de sólidos) nos dois equipamentos (MCR102 e reômetro capilar) revela uma notável concordância, particularly na região onde as faixas de taxa de cisalhamento se sobrepõem. Tanto as curvas de fluxo (Figura \ref{fig:fluxo_comparativo}) quanto as curvas de viscosidade (Figura \ref{fig:viscosidade_comparativa}) mostram uma proximidade significativa entre os dois conjuntos de dados. Este resultado sustenta a confiabilidade das medições realizadas com o dispositivo de baixo custo, especialmente para taxas de cisalhamento mais elevadas (tipicamente acima de \SI{100}{\per\second}), faixa onde os reômetros capilares costumam operar com maior eficácia.
    \framebox{\parbox{0.8\textwidth}{\centering
        \vspace{3cm}
        \textbf{Figura 5 Placeholder} \\
        \small\textit{Comparativo das curvas de viscosidade aparente em função da taxa de cisalhamento em escala logarítmica.}
        \vspace{3cm}
      }}
    \caption{Comparativo das curvas de viscosidade aparente em função da taxa de cisalhamento em escala logarítmica, obtidas no reômetro comercial (MCR102) e no reômetro capilar de baixo custo.}
    \label{fig:viscosidade_comparativa}
\end{figure}

A inspeção dos gráficos indica que todas as amostras testadas exibem um comportamento reológico não-newtoniano do tipo pseudoplástico, também conhecido como \textit{shear-thinning}. Este comportamento é evidenciado pela redução da viscosidade aparente conforme a taxa de cisalhamento aumenta, sendo uma característica bem documentada para suspensões aquosas de caulim na literatura científica.

O achado central desta seção reside na validação do reômetro capilar de baixo custo. A comparação dos dados obtidos para a Amostra F (60\% de sólidos) nos dois equipamentos (MCR102 e reômetro capilar) revela uma notável concordância, particularmente na região onde as faixas de taxa de cisalhamento se sobrepõem. Tanto as curvas de fluxo (Figura \ref{fig:fluxo_comparativo}) quanto as curvas de viscosidade (Figura \ref{fig:viscosidade_comparativa}) mostram uma proximidade significativa entre os dois conjuntos de dados. Este resultado sustenta a confiabilidade das medições realizadas com o dispositivo de baixo custo, especialmente para taxas de cisalhamento mais elevadas (tipicamente acima de \SI{100}{\per\second}), faixa onde os reômetros capilares costumam operar com maior eficácia.

A Tabela \ref{tab:parametros_reologicos} apresenta os parâmetros ajustados do modelo de Herschel-Bulkley para as amostras analisadas no reômetro capilar. Destaca-se o efeito pronunciado da concentração de sólidos sobre a tensão de escoamento ($\tau_0$). Observa-se um aumento drástico deste parâmetro, partindo de um valor próximo a zero para a amostra com 60\% de sólidos e ultrapassando \SI{450}{\pascal} para a amostra com 65\% de sólidos. Este aumento expressivo na tensão mínima necessária para iniciar o fluxo oferece uma explicação quantitativa para a impossibilidade de testar as amostras mais concentradas (A e B). A força de escoamento requerida excedeu a capacidade de pressurização do sistema pneumático do reômetro capilar.

\begin{table}[h!]
\centering
\caption{Parâmetros reológicos do modelo de Herschel-Bulkley para as amostras analisadas no reômetro capilar de baixo custo.}
\label{tab:parametros_reologicos}
\begin{tabular}{@{}ccccc@{}}
\toprule
\textbf{Amostra} & \textbf{\% Sólidos} & \textbf{$\tau_0$ (\si{\pascal})} & \textbf{K (\si{\pascal\cdot s^n})} & \textbf{n (-)} \\ \midrule
F & 60,0\% & $\approx$ 0 & 77,58 & 0,22 \\
E & 62,5\% & 6,31 & 103,90 & 0,19 \\
D & 65,0\% & 451,35 & 0,09 & 1,02 \\ \bottomrule
\end{tabular}
\end{table}

Esta constatação possui implicações diretas para o desenvolvimento de pastas destinadas à técnica DIW. Em aplicações de impressão 3D, busca-se maximizar o teor de sólidos na pasta para minimizar a retração durante a secagem e, assim, obter peças finais mais densas e com maior resistência mecânica. A correlação entre baixa absorção de água (indicativo de baixa porosidade) e alta resistência mecânica em cerâmicas foi confirmada por \citeonline{thurler2000}. Diante do aumento exponencial da tensão de escoamento observado, torna-se claro que a formulação de pastas de caulim puro com alto teor de sólidos (>65\%) para impressão é impraticável sem o uso de aditivos. A incorporação de plastificantes, como a bentonita (cuja importância para a impressão de metacaulim foi evidenciada por \citeonline{campos2023}), ou de defloculantes, como o poliacrilato de sódio, representa uma estratégia necessária para modular a tensão de escoamento e a viscosidade, viabilizando a extrusão em altas concentrações de sólidos \citeonline{revelo2018, zocca2015}.

\section{Conclusões}
\label{sec:conclusoes}
% Texto reescrito
A investigação reológica das suspensões aquosas de caulim confirmou seu comportamento característico de afinamento por cisalhamento (\textit{shear-thinning}). Foi observado um aumento significativo da viscosidade e, mais notadamente, da tensão de escoamento, à medida que a concentração de partículas sólidas na suspensão foi elevada.

A contribuição central deste estudo reside na concepção, construção e validação de um reômetro capilar de baixo custo. A análise comparativa dos resultados obtidos com este dispositivo em relação a um reômetro rotacional comercial estabelecido como referência demonstrou uma boa correlação entre as medições. Este resultado valida o equipamento desenvolvido como uma alternativa acessível e confiável para a caracterização reológica de suspensões cerâmicas, especialmente em regimes de cisalhamento mais intensos. Adicionalmente, as próprias limitações encontradas na operação do reômetro de baixo custo (a incapacidade de processar amostras com teor de sólidos superior a 65\%) serviram como corroboração experimental do rápido aumento da consistência das pastas ao se elevar a concentração de sólidos.

Este trabalho ilustra o potencial do uso de ferramentas de baixo custo para a realização de análises de materiais com um nível de precisão adequado para muitos propósitos de pesquisa e desenvolvimento. As caracterizações reológicas realizadas e a metodologia experimental estabelecida fornecem uma base sólida para investigações futuras. Tais investigações poderão se concentrar na otimização de formulações de caulim por meio da adição estratégica de modificadores reológicos (plastificantes, defloculantes, cargas minerais), com o objetivo final de desenvolver pastas cerâmicas com propriedades otimizadas para aplicações em manufatura aditiva por extrusão (impressão 3D).


% --- BIBLIOGRAFIA ---
\bibliography{Referencias} % Usa o arquivo Referencias.bib

\end{document}


